\documentclass[aspectratio=169]{beamer}

\usetheme{metropolis}
\usepackage{appendixnumberbeamer}

\usepackage{amsmath, amssymb}
\usepackage{bm}
\usepackage{graphicx}
\usepackage{booktabs}
\usepackage{xcolor}

\definecolor{accent}{RGB}{0, 82, 136}
\setbeamercolor{frametitle}{bg=accent}
\setbeamercolor{progress bar}{fg=accent}
\setbeamercolor{alerted text}{fg=accent}
\metroset{block=fill, progressbar=frametitle}

\title[XBRL2Vec]{\textbf{XBRL2Vec}\\[4pt]
  \large Compact Embeddings for Company Fundamentals\\with Macroeconomic Context}
\author{N.\ Serraino \and P.\ Zuccolotto \and A.\ Locoro}
\institute{University of Brescia \quad$\cdot$\quad Intesa Sanpaolo}
\date{}

% ─────────────────────────────────────────────────────────────────────────────
\begin{document}

\begin{frame}
  \titlepage
\end{frame}

% ─────────────────────────────────────────────────────────────────────────────
\begin{frame}{The Analogy: Word2Vec for Companies}

  \begin{columns}[T]
    \begin{column}{0.48\textwidth}
      \textbf{What Word2Vec did for words}
      \begin{itemize}
        \item Maps each word to a dense vector in a learned space
        \item Similar words land close together: \textit{king} $\approx$ \textit{queen}
        \item Enables arithmetic: $\vec{\text{king}} - \vec{\text{man}} + \vec{\text{woman}} \approx \vec{\text{queen}}$
        \item One representation, many downstream tasks
      \end{itemize}
    \end{column}
    \begin{column}{0.48\textwidth}
      \textbf{What XBRL2Vec does for companies}
      \begin{itemize}
        \item Maps each company (at each quarter) to a dense vector in a learned space
        \item Financially similar firms land close together
        \item Captures both firm-level dynamics \emph{and} macro exposure
        \item One representation: forecasting, peer selection, risk monitoring, stress testing
      \end{itemize}
    \end{column}
  \end{columns}

  \vspace{0.8em}
  \begin{block}{The key difference from Word2Vec}
    The embedding is not static — it is \emph{conditioned on the macroeconomic
    environment}.  The same company encodes differently under different macro
    regimes, capturing context that raw accounting figures cannot.
  \end{block}

\end{frame}

% ─────────────────────────────────────────────────────────────────────────────
\begin{frame}{The Problem: Financial Data Is Structured — and Context-Dependent}

  \begin{columns}[T]
    \begin{column}{0.48\textwidth}
      \textbf{What we have}
      \begin{itemize}
        \item Quarterly XBRL filings from SEC (balance sheet, income statement, cash flow)
        \item Macroeconomic indicators (rates, inflation, GDP growth, \ldots)
        \item A need to \emph{represent} companies, not just predict single numbers
      \end{itemize}
    \end{column}
  \end{columns}

  \vspace{1em}
  \begin{block}{Core idea}
    Learn a compact latent representation of each company's financial trajectory
    that is \emph{conditioned on the macroeconomic environment} — useful for
    both forecasting and downstream analytics.
  \end{block}

\end{frame}

% ─────────────────────────────────────────────────────────────────────────────
\begin{frame}{Architecture at a Glance}

  \begin{columns}[T]
    \begin{column}{0.65\textwidth}
      \begin{enumerate}
        \item \textbf{DLinear Financial/Macro Encoder} — decomposes each accounting/macro item into trend + seasonal components and projects them to a fixed-size embedding $\mathbf{h}_{\mathrm{fin}} \in \mathbb{R}^{2F}$
        \item \textbf{FiLM Conditioning} — macro embedding rescales and shifts
          the financial embedding feature-wise:
          \[
            \mathbf{h}'_{\mathrm{fin}} = \bm{\gamma} \odot \mathbf{h}_{\mathrm{fin}} + \bm{\delta}
          \]
          where $\bm{\gamma}, \bm{\delta}$ are predicted from $\mathbf{h}_{\mathrm{mac}}$
      \end{enumerate}
      The conditioned embedding is then projected to a latent vector
      $\mathbf{z}$, and a linear decoder reconstructs the financial forecast.
    \end{column}
    \begin{column}{0.42\textwidth}
      \centering
      \includegraphics[width=0.7\textwidth]{els-cas-templates/figs/full-architecture.png}\\
      {\footnotesize \textit{End-to-end Architecture}}
    \end{column}
  \end{columns}

\end{frame}

% ─────────────────────────────────────────────────────────────────────────────
\begin{frame}{Macro Personalisation: Not One-Size-Fits-All}

  \begin{columns}[T]
    \begin{column}{0.55\textwidth}
      Different companies respond differently to the same macroeconomic shock
      (a rate hike matters more to a bank than to a tech firm).

      \vspace{0.8em}
      Before feeding macro data to the encoder, we \textbf{re-weight each
      macro indicator per company} using XGBoost feature importances trained
      to predict the company's aggregate financial signal.

      \vspace{0.8em}
      This produces a personalised macro vector for each firm that reflects
      its historical sensitivity to each macroeconomic driver.
    \end{column}
    \begin{column}{0.42\textwidth}
      \begin{block}{In practice}
        \begin{itemize}
          \item $M = 7$ macro indicators (rates, inflation, GDP, \ldots)
          \item One XGBoost model per company, fitted on training data
          \item Importance weights applied before encoding
        \end{itemize}
      \end{block}
      \vspace{0.5em}
      \textit{Result: the macro signal reaching the FiLM layer already
      encodes firm-level factor exposure.}
    \end{column}
  \end{columns}

\end{frame}

% ─────────────────────────────────────────────────────────────────────────────
\begin{frame}{Does Macro Context Actually Help? — Forecasting Results}

  \textbf{Setup:} 18-configuration grid (input window $\times$ forecast horizon $\times$ latent size).
  Each run trains \textbf{XBRL2Vec} alongside a macro-blind ablation (same architecture, no macro).

  \vspace{0.8em}
  \begin{columns}[T]
    \begin{column}{0.48\textwidth}
      \begin{block}{Key finding}
        The contextual model outperforms the blind baseline in \textbf{all 18
        configurations} on out-of-sample MSE\@.  The improvement is
        statistically significant ($p < 0.001$, Wilcoxon signed-rank test,
        effect size $r = 1.0$).
      \end{block}
      \vspace{0.5em}

    \end{column}
    \begin{column}{0.48\textwidth}
      \centering
      \small
      \begin{tabular}{lcc}
        \toprule
        \textbf{Model} & \textbf{OOS MSE (mean)} \\
        \midrule
        XBRL2Vec (contextual) & \textbf{0.0205} \\
        XBRL2Vec-Blind        & 0.0210 \\
        \midrule
        \multicolumn{2}{l}{\textit{Macro exposure (Integrated Gradients):}} \\
        Latent $\mathbf{z}$          & 42.8\% of attribution \\
        Reconstruction $\hat{\mathbf{x}}$ & 11.3\% of attribution \\
        \bottomrule
      \end{tabular}
    \end{column}
  \end{columns}

\end{frame}

% ─────────────────────────────────────────────────────────────────────────────
\begin{frame}{What Drives the Latent Representation? — Feature Attribution}

  Attribution scores via \textbf{Integrated Gradients} (target: latent norm
  $\|\mathbf{z}\|_2$) show which inputs most shape the embedding.

  \begin{center}
    \includegraphics[width=0.72\textwidth]{els-cas-templates/figs/SALIENCY - global_features_latent_tin20-tout4-latent80.png}
  \end{center}
  {\footnotesize Top-20 features by global attribution.  Macro variables
  consistently rank among the most influential inputs alongside core
  financial items — confirming that the FiLM conditioning is actively used.}

\end{frame}

% ─────────────────────────────────────────────────────────────────────────────
\begin{frame}{Macro Exposure Is Heterogeneous Across Sectors}

  Same model, different target: attribution towards reconstruction error
  reveals \emph{which sectors rely most on macro signals}.

  \begin{center}
    \includegraphics[width=0.65\textwidth]{els-cas-templates/figs/SALIENCY - sector_bubble_reconstruction_tin20-tout4-latent80.png}
  \end{center}
  {\footnotesize Bubble size = number of companies.  Banks and federally
  sponsored credit agencies show the highest macro sensitivity — recovered
  \emph{without} any explicit sector label during training.}

\end{frame}

% ─────────────────────────────────────────────────────────────────────────────
\begin{frame}{Beyond Forecasting: Embedding-Based Peer Selection}

  \begin{columns}[T]
    \begin{column}{0.52\textwidth}
      \centering
      \includegraphics[width=\textwidth]{els-cas-templates/figs/n_corr.png}\\
      {\footnotesize Mean Pearson $\rho$ of future 6-quarter trajectories for
      $k{=}5$ neighbours selected by embedding similarity, raw financial
      similarity, or at random.  The embedding advantage is consistent
      across all evaluation quarters.}
    \end{column}
    \begin{column}{0.45\textwidth}
      \vspace{1.5em}
      \small
      \begin{tabular}{lcc}
        \toprule
        \textbf{Criterion} & \textbf{Mean} $\rho$ & \textbf{Median} $\rho$ \\
        \midrule
        Embedding & \textbf{0.244} & \textbf{0.293} \\
        Financial & 0.187 & 0.204 \\
        Random    & 0.029 & 0.038 \\
        \bottomrule
      \end{tabular}
      \vspace{1em}

      \normalsize
      \begin{block}{Takeaway}
        Embedding neighbours show \textbf{30\% higher future co-movement}
        than accounting-space neighbours.  The latent space encodes
        forward-looking financial structure not visible in the raw filings.
      \end{block}
    \end{column}
  \end{columns}

\end{frame}

% ─────────────────────────────────────────────────────────────────────────────
\begin{frame}{Potential Applications}

  \begin{columns}[T]
    \begin{column}{0.48\textwidth}
      \begin{block}{Credit risk monitoring}
        Track shifts in a company's latent position over time.
        Drift towards distressed clusters can serve as an early-warning signal.
      \end{block}
      \vspace{0.5em}
      \begin{block}{Portfolio construction}
        Use embedding-space similarity as an economically grounded alternative
        to raw accounting ratios for peer grouping and diversification.
      \end{block}
    \end{column}
    \begin{column}{0.48\textwidth}
      \begin{block}{Macro stress testing}
        Inject hypothetical macro vectors through the FiLM conditioner and
        observe how company embeddings shift — without retraining.
      \end{block}
      \vspace{0.5em}
      \begin{block}{Regime detection}
        Monitor cross-sectional shifts in the latent distribution as a
        market-level diagnostic for structural economic changes.
      \end{block}
    \end{column}
  \end{columns}

\end{frame}

% ─────────────────────────────────────────────────────────────────────────────
\begin{frame}{Summary}

  \begin{columns}[T]
    \begin{column}{0.55\textwidth}
      \textbf{What XBRL2Vec delivers:}
      \begin{itemize}
        \item A compact, macro-conditioned embedding of company fundamentals
          built entirely on open, standardised XBRL filings
        \item Improved out-of-sample forecasting across all tested configurations
        \item A latent space that is geometrically faithful, economically
          interpretable, and analytically reusable
        \item Firm-level macro exposure estimates without explicit sector labels
      \end{itemize}
    \end{column}
    \begin{column}{0.42\textwidth}
      \textbf{Next steps:}
      \begin{itemize}
        \item Richer macro feature set (credit spreads, sector indices)
        \item Feature-level macro weighting ($M \times F$ importance matrix)
        \item Extension to private companies and non-US jurisdictions
        \item Integration with downstream risk and valuation workflows
      \end{itemize}
      \vspace{1em}
      \begin{block}{}
        \centering Code and data pipeline will be released open-source.
      \end{block}
    \end{column}
  \end{columns}

\end{frame}

\end{document}
